\documentclass[12pt]{article}
\usepackage[margin=1in]{geometry}
\usepackage{amsmath}
\usepackage{fancyhdr}
\fancypagestyle{plain}{%
  \fancyhf{}
  \fancyhead[L]{POLISCI 450A}
  \fancyhead[R]{Fall 2026}
  %\fancyfoot[C]{\thepage}
}

\usepackage{titlesec}
\titleformat{\section}{\large\bfseries}{\thesection}{1em}{}
\usepackage{titling}
\setlength{\droptitle}{-3em}   % move title up


\title{\Large Lecture 2-3: Hypothesis Testing}
\author{}
\date{}
\begin{document}
\maketitle
\vspace{-6em}

\noindent \textbf{Goal:} Understand how to conduct hypothesis testing, interpret $p$-values, and their link to confidence intervals.\medskip

\begin{itemize}\setlength{\itemsep}{1em}
\item What is hypothesis testing?
\begin{itemize}
\item Trial analogy
\item Type I and II errors
\item Significance level $\alpha$
\item Power and $1-\beta$
\end{itemize}
\item Steps of hypothesis testing
\begin{itemize}
\item State $H_0$ and $H_1$
\item Choose $\alpha$
\item Select test statistic
\item Derive null distribution
\item Use critical values
\item Decision rule
\end{itemize}
\item Hypothesis testing with a single variable
\begin{itemize}
\item One-sample $t$-test setup
\item Standardized mean statistic
\item Null distribution: $t$ or normal
\item One- vs two-sided tests
\item Critical values and rejection regions
\end{itemize}
\item $p$-values and confidence intervals
\begin{itemize}
\item Definition of $p$-value
\item Direction depends on alternative
\item Invert tests to intervals
\item Test via confidence interval inclusion
\end{itemize}
\item Practical vs statistical significance
\begin{itemize}
\item Effect size vs precision
\item Large $n$ inflates significance
\item Report substantive importance
\end{itemize}
\end{itemize}


% End outline
\clearpage


\section*{Worksheet}

\begin{enumerate}\itemsep1.2em

% \item (\emph{Size control from critical values})
% Consider the two–sided one–sample $t$ test of $H_0:\mu=\mu_0$ that rejects when $\lvert T\rvert>c$, where $T=\dfrac{\overline{X}-\mu_0}{S/\sqrt{n}}$ and, under the normal model, $T\sim t_{n-1}$ when $H_0$ holds. Show that choosing $c=t_{1-\alpha/2}(n-1)$ yields a level–$\alpha$ test, i.e.\ $\Pr(\text{reject }H_0\mid H_0)=\alpha$.

% {\red Under $H_0$ and normality, $T\sim t_{n-1}$ and the $t$ distribution is symmetric. Hence

% $$
% \Pr(\text{reject }H_0\mid H_0)=\Pr(|T|>c)=2\,\Pr(T>c).
% $$

% Setting $c=t_{1-\alpha/2}(n-1)$ gives $\Pr(T>c)=\alpha/2$, so $\Pr(|T|>c)=2(\alpha/2)=\alpha$.}

% \item (\emph{$p$–value formula for a two–sided $t$ test})
% Show that for an observed $t$ statistic $t_{\text{obs}}$, the two–sided $p$–value is

% $$
% p \;=\; 2\Big(1-F_{t_{n-1}}\big(|t_{\text{obs}}|\big)\Big),
% $$

% where $F_{t_{n-1}}$ is the CDF of $t_{n-1}$.

% {\red By definition, the two–sided $p$–value is the null probability of obtaining a value at least as extreme as $t_{\text{obs}}$ in both tails:

% $$
% p \;=\; \Pr\big(|T|\ge |t_{\text{obs}}|\mid H_0\big)
% = \Pr\big(T\le -|t_{\text{obs}}|\big)+\Pr\big(T\ge |t_{\text{obs}}|\big).
% $$

% Symmetry of $t_{n-1}$ implies $\Pr(T\le -a)=\Pr(T\ge a)=1-F_{t_{n-1}}(a)$ for $a\ge0$, yielding

% $$
% p \;=\; 2\big(1-F_{t_{n-1}}(|t_{\text{obs}}|)\big).
% $$}

\item (\emph{Test–CI equivalence by inversion})  
For the two–sided level–$\alpha$ $t$ test of $H_0:\mu=\mu_0$, prove that $H_0$ is rejected if and only if $\mu_0$ is not contained in the $100(1-\alpha)\%$ $t$–confidence interval
\[
\left[\;\overline{X}-t_{\alpha/2}(n-1)\frac{S}{\sqrt{n}}\;,\;\overline{X}+t_{\alpha/2}(n-1)\frac{S}{\sqrt{n}}\;\right].
\]
\vspace{16em}

\item (\emph{Invariance to unit changes})
Let $Y_i=a+bX_i$ with $b\neq0$. Consider testing $H_0:\mu=\mu_0$ using the $t$ statistic based on ${X_i}$, and testing $H_0':\mu_Y=a+b\mu_0$ using the $t$ statistic based on ${Y_i}$. Show that both two–sided tests yield the same rejection decision.

\clearpage


\item (\emph{Power depends on effect size, variance, and $n$})
Assume known $\sigma$ and test $H_0:\mu=\mu_0$ vs.\ two–sided $H_1$ using the $z$ statistic $Z=\dfrac{\overline{X}-\mu_0}{\sigma/\sqrt{n}}$, rejecting for $|Z|>z_{\alpha/2}$. Derive the power function at $\mu=\mu_0+\delta$ and show: (i) for any fixed $\delta\neq0$, power $\to1$ as $n\to\infty$; (ii) for fixed $n$, power increases in $|\delta|$ and decreases in $\sigma$. (optional)

\end{enumerate}


\end{document}
